%% ma: L12-B02-0011
%% lop: 12
%% chuong: 1
%% bai: 2
%% dang: 12.2.6
%% loai: TN4
%% muc: TH
%% dap_an: "C"
%% nguon: "(NBV)-ĐỀ SỐ 13-MỖI NGÀY 1 ĐỀ THI 2025.pdf (D13, MỖI NGÀY 1 ĐỀ - Đề số 13), Phần I Câu 1"
%% kien_thuc: "Bài 2 (giá trị lớn nhất, nhỏ nhất trên một đoạn đọc từ đồ thị)"
%% trang_thai: chinh-thuc
%% ghi_chu: "Đáp án gốc C đúng (M = 6 tại x = 1, m = 1 tại x = -5). Hình TikZ dựng lại từ số liệu đọc trên hình gốc (các điểm cực trị và đầu mút), đường cong chỉ mang tính minh họa."
\begin{ex}
Cho hàm số $y=f(x)$ liên tục trên đoạn $[-5;3]$ và có đồ thị như hình vẽ. Gọi $M$ và $m$ lần lượt là giá trị lớn nhất và giá trị nhỏ nhất của hàm số trên đoạn $[-5;3]$. Giá trị $M+m$ bằng
\begin{center}
\begin{tikzpicture}[x=.62cm,y=.5cm]
  \draw[phu] (-5,0)--(-5,1)--(0,1) (-4,0)--(-4,3)--(0,3) (-3,0)--(-3,2)--(0,2) (-2,0)--(-2,4)--(0,4) (-1,0)--(-1,2) (1,0)--(1,6)--(0,6) (3,0)--(3,2)--(-1,2);
  \draw[truc] (-5.6,0)--(3.9,0) node[below]{$x$};
  \draw[truc] (0,-.3)--(0,7) node[right]{$y$};
  \draw[thick,blue!70!black] (-5,1) .. controls (-4.8,2.6) and (-4.4,3) .. (-4,3) .. controls (-3.5,3) and (-3.1,2.6) .. (-3,2)
    .. controls (-2.9,3.6) and (-2.5,4) .. (-2,4) .. controls (-1.5,4) and (-1.1,3.6) .. (-1,2)
    .. controls (-.9,5) and (.4,6) .. (1,6) .. controls (2,6) and (3,3.6) .. (3,2);
  \foreach \t in {-5,-4,-3,-2,-1,1,3} \draw (\t,.1)--(\t,-.1) node[below,font=\footnotesize]{$\t$};
  \foreach \t in {1,2,3,4,6} \draw (.07,\t)--(-.07,\t) node[left,font=\footnotesize]{$\t$};
  \node[below right,font=\footnotesize] at (0,0) {$O$};
  \node[font=\footnotesize] at (-3.7,5.2) {$y=f(x)$};
\end{tikzpicture}
\end{center}
\choice{$-2$}{$2$}{\True $7$}{$8$}
\loigiai{Trên đoạn $[-5;3]$, điểm cao nhất của đồ thị có tung độ $6$ (tại $x=1$) và điểm thấp nhất có tung độ $1$ (tại $x=-5$). Suy ra $M=6$, $m=1$ và $M+m=7$. Chọn C.}
\end{ex}
